\section{Introduction}\label{sec:introduction}
A neural solution of a dynamic economic model should be judged by the policy it delivers and the accuracy that can be established for that policy. Average residuals are not a substitute for such an assessment. Stopping boundaries make the distinction particularly important: a boundary observation may describe immediate liquidation without being the limit of continuation value from the interior. A numerical approximation that ignores this distinction can be inconsistent before the choice of optimizer or network width is considered.

This paper develops neural Bellman operators with explicit policy evaluation, improvement, and verification stages. The economic objective is an accurately certified state-contingent policy for the original continuous-time decision problem. We retain the full control opportunity and the original stopping contract when evaluating a candidate. The verification stage is independent of the optimizer's stopping rule and returns an a posteriori policy-loss bound. We also give a quantitative policy-iteration recurrence whose hypotheses are certified evaluation and improvement defects, rather than sampled training losses.

Our first construction concerns a stopped consumption--investment economy with costly preference adjustment. We prove a strictly positive jump, exceeding $6.37$ utility units, between continuation value near its upper wealth boundary and the liquidation payoff at that boundary. Thus a continuous critic constrained to equal the liquidation payoff on every face cannot acquire a vanishing uniform positive Bellman residual. This is a property of the unchanged economy, not a reason to alter its settlement. A signed, accessibility-weighted comparison resolves the inconsistency: the candidate policy cannot reach the upper wealth face, while its witness dominates settlement there and matches settlement on the reachable faces. These properties eliminate the entire trace contribution to its policy certificate without eliminating any action from the optimal upper comparison.

The associated full-horizon neural experiment retains ten seeds, two widths, and both prescribed checkpoints. Each stored policy and critic is enclosed over the complete original state--time cylinder. The boundary construction succeeds exactly, but the current residual bounds remain large: the twenty time-zero policy-loss bounds range from \RNNMin\ to \RNNMax. None meets $10^{-2}$. Six newly generated state-space candidates, from controlled Markov-chain and semi-Lagrangian schemes, are evaluated under the same continuous-time contract and common witness. They also do not meet this target. These calculations distinguish an effective correction to the approximation space from a claim that the current nonlinear training procedure is already competitive at high verified accuracy.

Our second application provides an accurate neural dynamic-control calculation in a coupled inventory economy with up to 128 state variables. Neural time-dependent gains are trained against differential Riccati residuals, without optimal-solution labels. A quadratic performance-difference identity and a Lyapunov moment bound turn a complete time enclosure into a certificate on the full state space. Every final checkpoint has per-coordinate loss bounded by \RInventoryFinalMax\ for initial mean squared inventory at most one. The method exploits two invariant modes; a structure-aware classical method can exploit exactly the same structure. The result establishes a fully certified high-dimensional neural policy, not a generic dimension-independent nonlinear solver or a neural advantage over scalar Riccati integration.

Sharp economic comparisons in the original nonlinear economy remain separate from these neural policies. A complete MPFR-directed arithmetic recalculation of all eighteen historical policy/dual nodes yields a price-uniform bound of \RMPFRGap. A fresh model-to-library experiment, starting only from model constants and deterministic optimization initializations, generates \RFreshNodes\ nodes and a bound of \RFreshGap\ over all adjustment prices in $[.5,8]$. The fresh experiment includes policy construction, dual fitting, adaptive node selection, original stopping evaluation, and exact rational price-envelope construction. We interpret its loss through a precisely specified, externally financed consumption-top-up experiment rather than calling an absolute utility tolerance an economic normalization.

These results have different domains and numerical roles. The original neural and classical feedback certificates cover the full nonlinear state--time domain but are not yet sharp. The time-control libraries give sharp original-economy comparisons at specified initial states and prices. The inventory construction gives a sharp state-global neural policy certificate in a different, structured economy. None of these statements substitutes one numerical object for another. The previous coefficient-transport, state--price continuation, adverse seed comparisons, external holdouts, and frozen-jet action-oracle results remain in the paper and supplement as identified historical evidence.

\subsection{Relation to the literature}
Stopped verification, contraction arguments, and the policy-improvement principle are classical. \citet{Judd1998} develops numerical dynamic programming and economically meaningful approximation assessment; \citet{KushnerDupuis2001} studies controlled Markov-chain approximation, and \citet{BarlesSouganidis1991} gives a fundamental monotone-approximation convergence framework. Adaptive sparse grids are an important computational-economics alternative \citep{BrummScheidegger2017}. Information and control relaxations underpin a substantial literature on dual upper bounds \citep{BrownSmithSun2010,BrownSmith2011}. Neural approaches include \citet{HanE2016,HanJentzenE2018}, and the retained external comparison uses the implementation of \citet{CaiFangZhou2025}.

The methodological content here is the identification and correction of a stopping-trace inconsistency in the neural approximation space, the execution of signed policy-specific boundary verification, the separation of certified policy-iteration defects from discretization errors, and a structure-exploiting state-global neural certificate. The affine-preference dual and joint state--price construction are specialized economic comparison tools. MPFR supplies established directed arithmetic \citep{FousseEtAl2007}; it does not independently prove the shared analytic quadrature or stopping lemmas. We make these dependencies explicit so that reproducibility, arithmetic independence, convergence, and demonstrated policy accuracy are not conflated.
